% --- Chapter: Network Programming ---
\section{Network Programming}

CRISP provides network I/O through high-level convenience functions for the most common tasks: fetching web resources, sending data to servers, and working with URLs.

\subsection{Fetching Web Resources}

The \texttt{http\_get} function fetches a URL and returns the response body as a string:

\begin{tcolorbox}[colback=codebg, colframe=darkpurple, colbacktitle=darkpurple, title=\textbf{\textcolor{codegold}{HTTP GET}}, fonttitle=\bfseries, sharp corners]
\begin{lstlisting}
# Fetch a web page
let html = http_get("https://example.com");
say "Fetched " + html.len() + " bytes";

# Fetch JSON from an API
let response = http_get("https://api.github.com/repos/user/crisp");
if json_valid(response) {
    let data = from_json(response);
    say "Stars: " + data.stargazers_count;
    say "Forks: " + data.forks_count;
}
\end{lstlisting}
\end{tcolorbox}

\subsection{Sending Data}

\texttt{http\_post} sends data to a URL with a specified content type:

\begin{tcolorbox}[colback=codebg, colframe=darkpurple, colbacktitle=darkpurple, title=\textbf{\textcolor{codegold}{HTTP POST}}, fonttitle=\bfseries, sharp corners]
\begin{lstlisting}
# Send JSON data
let payload = to_json({
    name => "CRISP",
    version => "0.1.5"
});

let response = http_post("https://httpbin.org/post",
                         payload,
                         "application/json");

say "Server responded: " + response;
\end{lstlisting}
\end{tcolorbox}

\subsection{URL Encoding and Decoding}

\begin{tcolorbox}[colback=codebg, colframe=darkpurple, colbacktitle=darkpurple, title=\textbf{\textcolor{codegold}{URL Encoding}}, fonttitle=\bfseries, sharp corners]
\begin{lstlisting}
# Encode special characters for URLs
say url_encode("hello world");      # hello%20world
say url_encode("name=CRISP&v=1");   # name%3DCRISP%26v%3D1

# Decode URL-encoded strings
say url_decode("hello%20world");    # hello world
say url_decode("name%3DCRISP");     # name=CRISP
\end{lstlisting}
\end{tcolorbox}

\subsection{A Complete Example: Simple Web Client}

\begin{tcolorbox}[colback=codebg, colframe=darkpurple, colbacktitle=darkpurple, title=\textbf{\textcolor{codegold}{Web Client (web\_client.csp)}}, fonttitle=\bfseries, sharp corners]
\begin{lstlisting}
# web_client.csp — Fetch and display web resources

fn fetch_json_api(url) {
    say "Fetching: " + url;
    let raw = http_get(url);

    if raw == "" {
        say "Error: no response from server.";
        return null;
    }

    if !json_valid(raw) {
        say "Error: response is not valid JSON.";
        return null;
    }

    return from_json(raw);
}

# Try fetching a public API
let data = fetch_json_api("https://api.github.com/repos/rust-lang/rust");

if data != null {
    say "Repository: " + data.full_name;
    say "Description: " + data.description;
    say "Stars: " + data.stargazers_count;
    say "Language: " + data.language;
}
\end{lstlisting}
\end{tcolorbox}

\subsection{Network Function Reference}

\begin{table}[H]
\centering
\rowcolors{2}{softviolet!30}{white}
\caption{\textcolor{darkpurple}{Network Functions}}
\vspace{0.3em}
\begin{tabular}{@{}llp{8cm}@{}}
\toprule
\rowcolor{softvioletdark}
\textbf{\textcolor{white}{Function}} & \textbf{\textcolor{white}{Signature}} & \textbf{\textcolor{white}{Description}} \\
\midrule
\texttt{http\_get} & \texttt{http\_get(url) → String} & Fetch a URL via GET, return response body \\
\texttt{http\_post} & \texttt{http\_post(url, body, content\_type?) → String} & POST data to a URL \\
\texttt{url\_encode} & \texttt{url\_encode(s) → String} & Percent-encode special characters \\
\texttt{url\_decode} & \texttt{url\_decode(s) → String} & Decode percent-encoded string \\
\bottomrule
\end{tabular}
\end{table}

\vspace{1.5cm}
\begin{tcolorbox}[colback=lightlavender, colframe=softvioletdark, title=\textbf{\textcolor{darkpurple}{Design Note: HTTP Support}}, fonttitle=\bfseries, sharp corners]
CRISP's network module provides simple, blocking HTTP operations via the \texttt{ureq} crate—a minimal, dependency-light HTTP client. The design philosophy is:

\begin{itemize}
    \item \textbf{Simple by default:} \texttt{http\_get(url)} does one thing: fetches a URL and returns the body. No callbacks, no streams, no complexity.
    \item \textbf{Opt-in complexity:} Future releases will add streaming responses, async I/O, and WebSocket support for advanced use cases.
    \item \textbf{Self-contained:} The network functions are available globally (no \texttt{use network;} required) because web access is a fundamental scripting need, not a specialized system concern.
\end{itemize}

For low-level socket programming, CRISP delegates to the POSIX module (\texttt{use posix;} → \texttt{posix["socket"]}, \texttt{posix["connect"]}, etc.), giving you full control when needed.
\end{tcolorbox}
